\chapter{Conclusions and future work}
\label{chap:conclusions}

Report~2 asked whether third-order finite-volume methods provide enough additional accuracy
to justify their cost on serial, OpenMP, and CUDA hardware.  The answer is conditional but
clear.  Both WENO-3/SSPRK3 and ADER-3 recover third order on smooth one- and two-dimensional
advection, while WENO-2 provides the expected second-order reference.  HLLC preserves the
contact wave and, in the matched smooth comparison, reduces the WENO-3 and ADER-3 density
errors by about two thirds relative to the older HLL route.  Serial, OpenMP, and CUDA fields
agree to roundoff or to small localized differences associated with nonlinear fallback
decisions, and every accepted high-resolution discontinuous field retains positive density
and pressure.

The decisive result is the interaction between method and backend.  On the smooth
\(600^2\) problem, CUDA WENO-3 achieves \(L_1(\rho)=1.726\times10^{-8}\) in 2.71 s,
compared with \(4.703\times10^{-7}\) in 2.08 s for WENO-2.  ADER-3 gives a similar error,
\(1.681\times10^{-8}\), but requires 10.51 s.  WENO-3/HLLC on CUDA is therefore the best
overall accuracy--efficiency combination tested.  WENO-2 remains the appropriate economical
baseline, especially for small grids.  ADER-3 is numerically successful and exposes much
greater parallelism than either WENO route, but the present local predictor and robust
fallback make it too expensive to be the default choice.

The discontinuous cases reinforce rather than overturn this ranking.  WENO-3 retains
contacts more sharply than WENO-2 at a manageable CUDA cost.  ADER-3 captures the same wave
pattern and remains admissible, but its local sensor transitions are more visible and its
absolute runtime is highest.  Higher formal order is therefore not sufficient by itself:
the useful method is the one whose robustness mechanism, flux, and hardware implementation
preserve the accuracy benefit at an acceptable cost.

The phase-resolved measurements identify concrete next steps.  First, OpenMP work arrays
should be allocated once and reused, and adjacent loops should be fused where this reduces
barriers and memory traffic.  Second, CUDA CFL reduction, conservative update, and validity
checks should be fused or reorganised to reduce the residual synchronization exposed after
the spatial operator is accelerated.  Third, the ADER predictor and troubled-interface
criterion should be profiled together. A smoother or characteristic-based sensor may reduce
branch sensitivity without sacrificing positivity.  Characteristic-variable WENO and a
more selective fallback are also natural numerical extensions.

Further performance work should use reserved or otherwise controlled GPU measurements,
hardware counters, and a profiler such as Nsight to separate cache, bandwidth, occupancy,
and allocation effects.  Adaptive mesh refinement would then be a more informative route
to larger effective resolution than uniform grids beyond \(5000^2\).  A longer project
could additionally examine multi-GPU decomposition and vary the shock--bubble size to test
how active-region geometry affects thread occupancy and load balance.

The principal outcome is consequently not a universal winner independent of context.
It is a tested selection rule: use WENO-2 when low cost is primary, use WENO-3/HLLC on CUDA
when accuracy and robustness must both be retained, and treat ADER-3 as a promising one-step
route whose numerical machinery is ready but whose implementation still offers substantial
scope for optimization.
